\documentclass[12pt,a4paper]{article}
\usepackage{amsmath, amssymb}
\usepackage[margin=2cm]{geometry}
\usepackage{enumitem}

\begin{document}

\begin{center}
  \underline{\textbf{QUIZ 2 Genap 2023/2024}}
\end{center}

\vspace{0.5cm}

\begin{center}
  \begin{tabular}{ll}
    Mata Kuliah   & : Teori Peluang / 3 SKS       \\
    Hari, Tanggal & : Kamis, 13 Juni 2024         \\
    Waktu         & : 09.00-10.40 WIB (100 menit) \\
    Kelas         & : C
  \end{tabular}
\end{center}

\vspace{0.5cm}

\begin{enumerate}
  \item Jika $X$ dan $Y$ adalah variabel acak diskrit dengan pdf gabungan
        \[ f_{X,Y}(x,y) = \frac{c5^{x+y}}{x!y!}, \quad x=0,1,2,\cdots, y=0,1,2,\cdots, \]
        dan nol sebaliknya.
        \begin{enumerate}[label=(\alph*)]
          \item Temukan konstanta $c$.
          \item Temukan pdf marginal dari $X$ dan $Y$.
          \item Apakah $X$ dan $Y$ saling independen? Jelaskan alasan saudara.
        \end{enumerate}
        \textbf{Catatan:} Ingat bahwa $e^m = \sum_{z=0}^{\infty} \frac{m^z}{z!}$; untuk variabel acak diskrit $\mathop{\sum\!\sum}\limits_{(x,y)\in T_{X,Y}} f(x,y) = 1$, $f_X(x) = \sum_{y\in T_Y} f_{X,Y}(x,y)$

  \item Misalkan pdf bersama dari variabel acak $X$ dan $Y$ diberikan oleh
        \[ f_{X,Y}(x,y) = e^{-(x+y)} \quad 0 < x < \infty; 0 < y < \infty \]
        dan nol lainnya. Dapatkan:
        \begin{enumerate}[label=(\alph*)]
          \item CDF bersama $F(1,2)$;
          \item $P[X < Y]$.
        \end{enumerate}
        \textbf{Catatan:} $F_{X,Y}(x,y) = P(X \le x, Y \le y)$

  \item Misalkan $X$ dan $Y$ memiliki pdf gabungan
        \[ f_{X,Y}(x,y) = \frac{2}{3}(x+1), \quad 0 < x < 1, 0 < y < 1 \]
        dan nol sebaliknya.
        \begin{enumerate}[label=(\alph*)]
          \item Apakah $X$ dan $Y$ independen? Jelaskan alasan saudara.
          \item Dapatkan $P[X < Y < 2X]$.
        \end{enumerate}

  \item Misalkan $X$ dan $Y$ adalah variabel acak kontinu dengan pdf bersama,
        \[ f_{X,Y}(x,y) = k(x+y); \quad 0 \le x \le y \le 1 \]
        Dapatkan:
        \begin{enumerate}[label=(\alph*)]
          \item nilai $k$ sehingga $f_{X,Y}(x,y)$ adalah pdf bersama.
          \item pdf marginal, $f_X(x)$
          \item pdf bersyarat $f_{Y|X}(y|x)$
        \end{enumerate}
        \textbf{Catatan:} $\iint_{(x,y)\in T_{X,Y}} f(x,y) dx dy = 1$, $f_X(x) = \int_{T_Y} f_{X,Y}(x,y)dy$, $f_{Y|X}(y|x) = \frac{f_{X,Y}(x,y)}{f_X(x)}$
\end{enumerate}

\vspace{1cm}
\begin{center}
  \rule{3cm}{1.5pt} Good Luck \rule{3cm}{1.5pt}
\end{center}

\newpage
\section*{Solusi}

\begin{enumerate}
  \item
        \begin{enumerate}[label=(\alph*)]
          \item Syarat agar $f_{X,Y}(x,y)$ merupakan pdf (atau lebih tepatnya pmf untuk diskrit) adalah jumlah semua probabilitasnya sama dengan 1.
                \[ \sum_{x=0}^{\infty} \sum_{y=0}^{\infty} f_{X,Y}(x,y) = 1 \]
                \[ \sum_{x=0}^{\infty} \sum_{y=0}^{\infty} \frac{c5^{x+y}}{x!y!} = 1 \]
                Faktorkan ekspresi ke suku-suku yang memuat $x$ dan $y$:
                \[ c \left( \sum_{x=0}^{\infty} \frac{5^x}{x!} \right) \left( \sum_{y=0}^{\infty} \frac{5^y}{y!} \right) = 1 \]
                Berdasarkan sifat deret Maclaurin, $e^m = \sum_{z=0}^{\infty} \frac{m^z}{z!}$, sehingga:
                \[ c (e^5) (e^5) = 1 \]
                \[ c e^{10} = 1 \implies c = e^{-10} \]

          \item Pdf marginal dari $X$:
                \[ f_X(x) = \sum_{y=0}^{\infty} f_{X,Y}(x,y) = \sum_{y=0}^{\infty} \frac{e^{-10} 5^{x+y}}{x!y!} = \frac{e^{-10} 5^x}{x!} \sum_{y=0}^{\infty} \frac{5^y}{y!} = \frac{e^{-10} 5^x}{x!} (e^5) = \frac{e^{-5} 5^x}{x!} \]
                untuk $x = 0, 1, 2, \dots$

                Pdf marginal dari $Y$:
                \[ f_Y(y) = \sum_{x=0}^{\infty} f_{X,Y}(x,y) = \sum_{x=0}^{\infty} \frac{e^{-10} 5^{x+y}}{x!y!} = \frac{e^{-10} 5^y}{y!} \sum_{x=0}^{\infty} \frac{5^x}{x!} = \frac{e^{-10} 5^y}{y!} (e^5) = \frac{e^{-5} 5^y}{y!} \]
                untuk $y = 0, 1, 2, \dots$

          \item Ya, $X$ dan $Y$ saling independen. \\
                Alasan: Syarat independensi adalah $f_{X,Y}(x,y) = f_X(x) f_Y(y)$.
                \[ f_X(x) f_Y(y) = \left( \frac{e^{-5} 5^x}{x!} \right) \left( \frac{e^{-5} 5^y}{y!} \right) = \frac{e^{-10} 5^{x+y}}{x!y!} = f_{X,Y}(x,y) \]
                Karena persamaan ini memenuhi semua nilai $x, y \ge 0$, maka terbukti independen.
        \end{enumerate}

  \item
        \begin{enumerate}[label=(\alph*)]
          \item CDF bersama $F(1,2) = P(X \le 1, Y \le 2)$:
                \[ F(1,2) = \int_{0}^{1} \int_{0}^{2} e^{-(x+y)} \,dy \,dx = \left( \int_{0}^{1} e^{-x} \,dx \right) \left( \int_{0}^{2} e^{-y} \,dy \right) \]
                \[ F(1,2) = \left[ -e^{-x} \right]_0^1 \left[ -e^{-y} \right]_0^2 = \left( 1 - e^{-1} \right) \left( 1 - e^{-2} \right) \]

          \item $P[X < Y]$:
                Area integrasinya adalah $0 < x < \infty$ dan $y > x$.
                \[ P[X < Y] = \int_{0}^{\infty} \int_{x}^{\infty} e^{-(x+y)} \,dy \,dx = \int_{0}^{\infty} e^{-x} \left( \int_{x}^{\infty} e^{-y} \,dy \right) \,dx \]
                \[ = \int_{0}^{\infty} e^{-x} \left[ -e^{-y} \right]_x^{\infty} \,dx = \int_{0}^{\infty} e^{-x} (0 - (-e^{-x})) \,dx \]
                \[ = \int_{0}^{\infty} e^{-2x} \,dx = \left[ -\frac{1}{2} e^{-2x} \right]_0^{\infty} = 0 - \left(-\frac{1}{2}\right) = \frac{1}{2} \]
        \end{enumerate}

  \item
        \begin{enumerate}[label=(\alph*)]
          \item Untuk mengecek independensi, cari pdf marginal masing-masing:
                \[ f_X(x) = \int_{0}^{1} \frac{2}{3}(x+1) \,dy = \frac{2}{3}(x+1) [y]_0^1 = \frac{2}{3}(x+1), \quad 0 < x < 1 \]
                \[ f_Y(y) = \int_{0}^{1} \frac{2}{3}(x+1) \,dx = \frac{2}{3} \left[ \frac{1}{2}x^2 + x \right]_0^1 = \frac{2}{3} \left( \frac{1}{2} + 1 \right) = 1, \quad 0 < y < 1 \]
                Cek apakah $f_{X,Y}(x,y) = f_X(x) f_Y(y)$:
                \[ f_X(x) f_Y(y) = \left( \frac{2}{3}(x+1) \right) (1) = \frac{2}{3}(x+1) = f_{X,Y}(x,y) \]
                Ya, $X$ dan $Y$ independen.

          \item $P[X < Y < 2X]$:
                Batasan integrasi ditentukan oleh $y > x$ dan $y < 2x$ pada kotak $0 < x < 1$ dan $0 < y < 1$.
                Perpotongan garis $y = 2x$ dengan batas $y = 1$ berada di $x = 1/2$. Sehingga daerah ini dibagi dua:
                Untuk nilai elemen $x \in (0, 1/2]$, batas $y$ dari $x$ ke $2x$.
                Untuk nilai elemen $x \in (1/2, 1)$, batas $y$ dari $x$ ke $1$.
                \[ P = \int_{0}^{1/2} \int_{x}^{2x} \frac{2}{3}(x+1) \,dy \,dx + \int_{1/2}^{1} \int_{x}^{1} \frac{2}{3}(x+1) \,dy \,dx \]
                Bagian 1:
                \[ \int_{0}^{1/2} \frac{2}{3}(x+1) (2x - x) \,dx = \int_{0}^{1/2} \frac{2}{3}(x^2 + x) \,dx = \frac{2}{3} \left[ \frac{x^3}{3} + \frac{x^2}{2} \right]_0^{1/2} \]
                \[ = \frac{2}{3} \left( \frac{1/8}{3} + \frac{1/4}{2} \right) = \frac{2}{3} \left( \frac{1}{24} + \frac{1}{8} \right) = \frac{2}{3} \left( \frac{4}{24} \right) = \frac{1}{9} \]
                Bagian 2:
                \[ \int_{1/2}^{1} \frac{2}{3}(x+1) (1 - x) \,dx = \int_{1/2}^{1} \frac{2}{3}(1 - x^2) \,dx = \frac{2}{3} \left[ x - \frac{x^3}{3} \right]_{1/2}^{1} \]
                \[ = \frac{2}{3} \left( \left(1 - \frac{1}{3}\right) - \left(\frac{1}{2} - \frac{1}{24}\right) \right) = \frac{2}{3} \left( \frac{2}{3} - \frac{11}{24} \right) = \frac{2}{3} \left( \frac{5}{24} \right) = \frac{5}{36} \]
                Total probabilitas:
                \[ \frac{1}{9} + \frac{5}{36} = \frac{4}{36} + \frac{5}{36} = \frac{9}{36} = \frac{1}{4} \]
        \end{enumerate}

  \item
        \begin{enumerate}[label=(\alph*)]
          \item Nilai $k$ sehingga $f_{X,Y}(x,y)$ adalah pdf:
                Domain yang diberikan adalah $0 \le x \le y \le 1$. Integralkan terhadap $dx$ dulu dari $0$ ke $y$, lalu $dy$ dari $0$ ke $1$.
                \[ \int_{0}^{1} \int_{0}^{y} k(x+y) \,dx \,dy = 1 \]
                \[ \int_{0}^{1} k \left[ \frac{1}{2}x^2 + xy \right]_0^y \,dy = \int_{0}^{1} k \left( \frac{1}{2}y^2 + y^2 \right) \,dy = \int_{0}^{1} k \left( \frac{3}{2}y^2 \right) \,dy \]
                \[ k \left[ \frac{3}{2} \cdot \frac{y^3}{3} \right]_0^1 = k \left[ \frac{y^3}{2} \right]_0^1 = \frac{k}{2} \]
                Karena integral harus bernilai 1, maka $\frac{k}{2} = 1 \implies k = 2$.

          \item pdf marginal $f_X(x)$:
                Integralkan fungsi probabilitas bersama terhadap semua kemungkinan nilai $y$. Untuk nilai $x$ yang diam, batasan $y$ adalah dari $x$ sampai $1$.
                \[ f_X(x) = \int_{x}^{1} 2(x+y) \,dy = 2 \left[ xy + \frac{1}{2}y^2 \right]_x^1 \]
                \[ = 2 \left( \left(x + \frac{1}{2}\right) - \left(x^2 + \frac{1}{2}x^2\right) \right) = 2 \left( x + \frac{1}{2} - \frac{3}{2}x^2 \right) \]
                \[ = 1 + 2x - 3x^2, \quad \text{untuk } 0 \le x \le 1 \]

          \item pdf bersyarat $f_{Y|X}(y|x)$:
                Berdasarkan definisi pdf bersyarat:
                \[ f_{Y|X}(y|x) = \frac{f_{X,Y}(x,y)}{f_X(x)} \]
                \[ f_{Y|X}(y|x) = \frac{2(x+y)}{1 + 2x - 3x^2}, \quad \text{untuk } x \le y \le 1 \]
        \end{enumerate}
\end{enumerate}
\end{document}